\subsection{QIF integration and matched-carrier protocol}\label{supp:qif-protocol}
The finite network uses the PING model and observables defined in
\S\ref{sec:qif-methods}. The following settings support the weak-drive and
field-induced gamma comparisons.

Euler integration uses $0.0025$~ms steps and $800$~ms settling, exceeding
four field-free relaxation times at $\bar\eta=0.85$ ($176$~ms).
Measurements cover whole beat cycles to avoid phase-coverage bias:
$44$ at $55$~Hz and $34$ at $42$~Hz ($800$ and $810$~ms).

We also run both conditions at $A_f=10$ per tone with the same measurement
settings. For $\bar\eta=0.85$, separate checks double the population size,
double the measurement window, or halve the integration step; further runs raise
$V_{\rm peak}$ to $200$ and $400$ at both time steps. An equal-RMS
mean-field control replaces the two tones by one $300$~Hz carrier with peak
amplitude $\sqrt{2}A_f$, using the same duration: both inputs have mean-square
amplitude $A_f^2$ over these windows. The carrier comparison starts from the
field-free equilibrium and uses $800$~ms settling and $800$~ms measurement,
with a $3200$~ms window to check spectral leakage at the larger amplitude.

A longer window reduces the single carrier's apparent $55$~Hz
depth from $0.011$ to $9.2\times10^{-5}$, while TI's stays at $0.30$--$0.31$;
the former is spectral leakage.
